\chapter{Verification and Validation Approach}
\ChapterLabel{VerificationApproach}
\label{Section:VerificationApproachResponsibility}
This chapter explains how the architecture and requirements will be evaluated and how
objective evidence differs from planned outcomes.


\section{Evaluation Baseline}
The candidate contains twenty requirements with mapped tests, verification
entries, and validation entries. All product tests remain unexecuted and
all evidence dispositions remain Pending. The detailed
\hyperref[Table:ProjectTests]{test definitions} are blocked outlines pending
accepted contracts, reference inputs, fixtures, and controlled product
configurations. A complete set of trace links does not establish complete
stakeholder coverage or a passing product.

\section{Requirement Verification}
Verification compares observed behavior or inspected artifacts with the
specified obligation. Planned methods include source inspection and builds
for the C++ constraints, reference-input and transfer tests for acquisition,
association and qualification tests, timing measurements, record inspection,
fault injection, and interface-contract checks. The authoritative methods
and required evidence remain in
\hyperref[Table:VerificationRegister]{the verification register}; the
procedures and expected results remain in
\hyperref[Table:ProjectTests]{the test register}.

Each execution requires a controlled requirement and implementation baseline,
hardware and tool versions, resolved contracts, known or calibrated fixtures,
and approved expected data. Retained records identify input, procedure,
expected and observed results, date, operator, and evidence location under
\ChapterRef{TestEvidence}. An expected result is not an observed result.

\section{Stakeholder Validation}
Validation evaluates whether the system supports the intended stakeholder
use. Operator scenarios address reading interpretation; acquisition scenarios
address continuity under processing load; field scenarios address platform
use; integration scenarios address contract interpretation and compatible
replacement. The authoritative scenarios and success criteria remain in
\hyperref[Table:ValidationRegister]{the validation register}. Michael
Stoddard coordinates participants and evidence; Dr. Darian Muresan retains
the recorded disposition authority. No stakeholder observation or acceptance
is claimed by this publication.

\section{Architecture Evaluation and Open Coverage}
The \hyperref[Table:QualityScenarioCoverage]{quality scenario coverage}
connects architecture concerns to candidate evidence and owned decisions.
The 10 ms local publication deadline, 100 ms local stale-state limit, and
1 ms timestamp-error limit are candidate allocations under
\RequirementRef{REQ-TIM-01}, \RequirementRef{REQ-SMP-02}, and
\RequirementRef{REQ-TIM-02}. Mission delivery rate, platform scope,
interface compatibility, uncertainty, and recovery criteria require the
\hyperref[Section:ReviewDecisions]{technical review decisions} before
objective acceptance can be assessed.

\section{Traceability and Evidence Disposition}
The \hyperref[Table:TraceabilityMatrix]{traceability matrix} connects needs
to requirements and verification; the
\hyperref[Table:SupportingTraceLinks]{supporting links} connect tests, use
cases, architectural decisions, and validation. These relationships remain
current when an obligation or design changes. Verification and validation
are recorded separately, and missing evidence remains explicit. Successful
PDF compilation and source-link checks concern the document, not product
behavior or stakeholder acceptance.
